% from ega4-9-fr.tex
\oldpage[IV-3]{88}

l'extension algébriquement close $k'$ de $k$ considérée~; nous dirons que c'est
le \emph{type primaire géométrique} de $\mathcal F$. Ces définitions posées,
l'énoncé que nous avons en vue est le suivant~:

\phantomsection\label{IV.9.8.9.1-fr}%
\textbf{(9.8.9.1)} Soient $f:X\to S$ un morphisme de présentation finie,
$\mathcal F$ un $\mathcal O_X$-Module quasi-cohérent de présentation finie.
Pour tout $s\in S$, soit $T(s)$ le type primaire géométrique de $\mathcal F_s$.
Alors la fonction $s\mapsto T(s)$ est localement constructible dans $S$.

Compte tenu des remarques qui précèdent, on est ramené comme d'ordinaire à
prouver que si $S$ est affine, noethérien et intègre de point générique $\eta$,
$T(s)$ est constante au voisinage de $\eta$. Raisonnant comme au début de
(9.7.7), on peut supposer que toutes les parties irréductibles de $X_\eta$ qui
interviennent sont géométriquement irréductibles, et alors la proposition
résulte de (9.5.1), (9.5.5), (9.8.3) et (9.8.6)~; nous laissons les détails au
lecteur. On pourrait généraliser en considérant plusieurs Modules cohérents et
en définissant leur «~squelette primaire simultané~», etc. La conclusion
générale de ce qui a été vu depuis le début de ce paragraphe est que pour toutes
les propriétés du type envisagé (et pour un $S$ irréductible) \emph{les
propriétés valables sur la «~fibre générique~» le sont encore sur toutes les
fibres voisines}.
\end{remark}

\subsection{Constructibilité des propriétés locales des fibres}
\label{subsection:IV.9.9-fr}

\begin{proposition}[9.9.1]
\label{IV.9.9.1-fr}
Soient $f:X\to S$ un morphisme localement de présentation finie, $Z$ une
partie localement constructible de $X$ telle que pour tout $s\in S$, $Z_s$ soit
fermé dans $X_s$, $\Phi$ une partie finie de $\mathbf Z\cup\{\pm\infty\}$.
Alors les parties suivantes de $X$ sont localement constructibles~:
\begin{enumerate}
\item[(i)] L'ensemble des $x\in X$ tels que
$\dim_x(Z_{f(x)})$ appartienne à $\Phi$.
\item[(ii)] L'ensemble des $x\in X$ tels que
$\operatorname{codim}_x(Z_{f(x)},X_{f(x)})$ appartienne à $\Phi$.
\item[(iii)] L'ensemble des $x\in X$ tels que l'anneau local
$\mathcal O_{X_{f(x)},x}$ soit équidimensionnel.
\end{enumerate}
\end{proposition}

On notera que les propriétés (i) et (ii) s'expriment encore en disant que les
fonctions $x\mapsto\dim_x(Z_{f(x)})$ et
$x\mapsto\operatorname{codim}_x(Z_{f(x)},X_{f(x)})$ sont
\emph{localement constructibles} dans $X$ ($\mathbf 0_{\mathrm{III}},9.3.1$).

Les questions étant locales sur $X$, on peut se borner au cas où
$S=\operatorname{Spec}(A)$ et $X=\operatorname{Spec}(B)$ sont affines et où
$f$ est un morphisme de présentation finie~; il existe alors un sous-anneau
$A_0$ de $A$ qui est une $\mathbf Z$-algèbre de type fini, un
$A_0$-préschéma de type fini $X_0$ et une partie constructible $Z_0$ de $X_0$
tels que $X=X_0\otimes_{A_0}A$ et $Z=h^{-1}(Z_0)$, où $h:X\to X_0$ est la
projection canonique ((8.9.1) et (8.3.11)). En outre, pour tout $s\in S$, si
$s_0$ est la projection de $s$ dans $S_0=\operatorname{Spec}(A_0)$, on a
\[
X_s=(X_0)_{s_0}\otimes_{k(s_0)}k(s),
\]
et $h_s$ est la projection $X_s\to(X_0)_{s_0}$, on a
$Z_s=h_s^{-1}((Z_0)_{s_0})$. Comme le morphisme $h_s$ est fidèlement plat et
quasi-compact, l'hypothèse que $Z_s$ est fermé dans $X_s$ entraîne que
$(Z_0)_{s_0}$ est fermé dans $(X_0)_{s_0}$ (2.3.12).

Cela étant, la transitivité des fibres (\textbf{I}, 3.6.4) et la prop. (4.2.7)
entraînent que l'ensemble des dimensions des composantes irréductibles de
$Z_s$ contenant $x$ est le même que l'ensemble des dimensions des composantes
irréductibles de $(Z_0)_{s_0}$ contenant $x_0=h_s(x)$. En particulier, on a
$\dim_x(Z_s)=\dim_{x_0}((Z_0)_{s_0})$. D'autre part, si $Z_s^{(\beta)}$ sont les
composantes irréductibles de $Z_s$ contenant $x$ et $X_s^{(\alpha)}$ les
composantes irréductibles de $X_s$

